% ===================================================================
%  Front-tracking and VOF for TWO-DIMENSIONAL two-phase flow
%  Newtonian and viscoelastic, on uniform and adaptive meshes
% ===================================================================
%
% PARA O OVERLEAF: suba esta pasta inteira (main.tex, preambulo.tex,
% secoes/, dados/, referencias.bib) e marque main.tex como documento
% principal.  Compilador: pdfLaTeX.  Nenhuma imagem binaria -- os graficos
% saem do pgfplots lendo dados/*.dat.
%
% ESTE E' O ARTIGO BIDIMENSIONAL.  O tridimensional e' o companheiro em
% ../3d/, com a mesma estrutura.  Os dois compartilham a tese -- mesma
% maquinaria para os dois metodos, de modo que a diferenca medida seja da
% REPRESENTACAO DA INTERFACE e nao do arranjo -- e diferem no que a dimensao
% permite medir:
%
%   em 2D   ha' um benchmark publicado com valores de referencia tabelados
%           (Hysing et al.), entao a comparacao e' contra numero externo;
%   em 3D   nao ha' equivalente, e a comparacao e' entre os metodos e contra
%           oraculos de forma fechada.
%
% ESTADO DE CADA SECAO.  Cada secao abre com \status{}; defina \nostatus para
% suprimi-los na submissao.  Eles existem para que um paragrafo de
% planejamento nunca seja lido como resultado.
%
%   1 Introduction           a escrever
%   2 Formulation            PRONTA: identica a' do 3D, restrita ao plano
%   3 Numerical method       PRONTA quanto ao que esta' implementado
%   4 Verification           PARCIAL: os oraculos 2D precisam ser inventariados
%   5 Newtonian 2D           FECHADA: os dois metodos a t=3, 12.000 passos,
%                            cinco alvos contra o publicado, e o achado do
%                            vies compartilhado (cinco desvios negativos)
%   6 Viscoelastic 2D        NAO MEDIDO: o VOF tem duas corridas curtas
%                            exploratorias (101 passos) e o front-tracking
%                            tem configuracao mas nenhuma corrida de producao
%   7 Cost and performance   NAO MEDIDO em 2D
%   8 Conclusions            depende das acima

\documentclass[preprint,review,12pt]{elsarticle}
\journal{Journal of Computational Physics}

% ---------------------------------------------------------------- preambulo
% Comum aos dois artigos (2D e 3D).  Mantido num arquivo so' para que uma
% mudanca de estilo nao precise ser feita duas vezes e divergir na segunda.
\usepackage[T1]{fontenc}
\usepackage[utf8]{inputenc}
\usepackage[english]{babel}
\usepackage{lmodern}
\usepackage{microtype}
\usepackage{amsmath,amssymb}
\usepackage{booktabs}
\usepackage{adjustbox}
\usepackage{siunitx}
\usepackage{graphicx}
\usepackage{pgfplots}
\usepackage{pgfplotstable}
\usepackage{tikz}
\usepackage{xcolor}
\usepackage{capt-of}
\pgfplotsset{compat=1.18}

% OS NUMEROS VIVEM EM dados/*.dat e sao lidos por pgfplotstable, nunca
% digitados no texto.  Regerar uma medida troca o .dat e o texto acompanha,
% sem chance de a tabela e a prosa discordarem.
\newcommand{\dat}[1]{dados/#1}

% \status{...}: o que a secao tem de respaldo medido.  Definir \nostatus antes
% do \begin{document} faz os blocos sumirem na versao de submissao.
\newcommand{\status}[1]{%
  \ifdefined\nostatus\else
    \begingroup
    \par\medskip\noindent
    \colorbox{yellow!18}{\begin{minipage}{0.97\linewidth}
      \small\textbf{Status of this section.} #1
    \end{minipage}}\par\medskip
    \endgroup
  \fi}

% \missing{...}: uma lacuna pontual, visivel de proposito.
\newcommand{\missing}[1]{\textcolor{red!70!black}{\textbf{[missing:} #1\textbf{]}}}

% \caveat NAO e' o mesmo que \missing.  \missing marca o que falta MEDIR;
% \caveat marca o que foi medido mas com uma ressalva que muda a leitura.
% Confundir os dois esconderia um resultado-com-reserva entre pendencias.
\newcommand{\caveat}[1]{\par\smallskip\noindent\textcolor{orange!60!black}{%
  \textbf{[caveat:} #1\textbf{]}}\par\smallskip}

\usepackage[hidelinks]{hyperref}

% Descomente para a versao de submissao, sem os blocos de estado:
% \def\nostatus{}

\emergencystretch=3em

\begin{document}

\begin{frontmatter}

\title{Front-tracking and VOF for two-dimensional two-phase flow:
       Newtonian and viscoelastic, on uniform and adaptive meshes}

\author[icmc]{Antonio Castelo Filho}
\ead{castelo@icmc.usp.br}
\author[icmc]{\missing{co-authors}}

\affiliation[icmc]{organization={Instituto de Ci\^encias Matem\'aticas e de
  Computa\c{c}\~ao, Universidade de S\~ao Paulo},
  city={S\~ao Carlos}, state={SP}, country={Brazil}}

\begin{abstract}
\missing{abstract --- to be written last, once the viscoelastic pair has been
run. What it has to say, if the
measurements hold: that the two interface methods share every component except
the interface representation itself, so that the measured difference is
attributable to that choice; that both reproduce the published rise-velocity,
circularity and centre-of-mass series of the Hysing benchmark in the
ellipsoidal regime, and how they differ in the skirted regime where the
reference solutions themselves disagree; that all five deviations from the
published values fall on the same side, in two independent discretisations,
which is the signature of a bias in what the methods share rather than in the
interface treatment; what viscoelasticity does to that comparison; and what adaptive refinement costs and returns in each --- including
the finding that a refinement criterion specified as a fixed \emph{cell count}
rather than a fixed \emph{physical band} makes a mesh-convergence study
impossible, because each resolution then refines a different physical width.}
\end{abstract}

\begin{keyword}
front-tracking \sep volume of fluid \sep viscoelastic two-phase flow \sep
rising bubble benchmark \sep adaptive mesh refinement \sep hierarchical grids
\end{keyword}

\end{frontmatter}

\tableofcontents

\section{Introduction}

\status{To be written. It has to do three things: place the two interface
families against the two-dimensional literature; state why a \emph{controlled}
comparison needs shared machinery rather than two codes; and say what the
viscoelastic case adds that the Newtonian one cannot show. The companion paper
\missing{cite the 3D companion} does the same in three dimensions, and the two
should be readable independently.}

\missing{introduction}

\section{Formulation}

\status{Ready. The equations are those of the companion three-dimensional
paper restricted to the plane; the solver's non-dimensionalisation convention
is the same and was verified against the code, not assumed from the
literature.}

\missing{formulation --- port from the 3D companion, restricted to the plane.
The four subsections carry over unchanged in content: governing equations; the
distinction between the polymeric tensor and the tensor the solver stores; the
constitutive model; and the solver's non-dimensionalisation convention. Only
the index range changes.}

\section{Numerical method}\label{sec:method}

\status{Ready as to what is implemented. The shared-machinery argument is the
same as in the three-dimensional companion and is verified in the code, not
assumed. The two-dimensional specifics below --- the normal and curvature
estimator, and the form of the front --- were read from the source, and where
the two dimensions diverge that is stated rather than smoothed over.}

\subsection{Shared machinery, and the single point of difference}

The comparison in this paper is only meaningful if the two methods differ in
\emph{one} thing. They are therefore built on the same components: the same
hierarchical grid, the same interpolated properties, the same momentum solver,
the same projection, and coupling at the same point of the multiphase step.
What differs is how the interface is represented, and everything the
representation entails --- how the normal and the curvature are obtained, and
how the interface is transported.

\missing{the itemised correspondence between the two paths through the
multiphase step, as in the 3D companion, with the line numbers of the single
branch point.}

\subsection{VOF: geometric PLIC in the plane}

The volume fraction is advected by a directionally split geometric scheme: one
sweep per direction, each reconstructing a straight segment per interface cell
from the fraction in its neighbourhood, fluxing the area that segment cuts, and
copying the auxiliary field back before the next sweep.

The normal and the curvature come from a \emph{height-function} estimator. The
two-dimensional path uses the variant of Shirani; the three-dimensional
companion uses the standard height function with Mehta's correction. This is a
genuine difference between the two papers and not a presentational one: the
estimator is part of what is being compared, and quoting a 2D number against a
3D one without saying so would be comparing two different methods.

\caveat{The code also carries ELVIRA and an adaptive height-function variant,
both reachable and both currently disabled at the call site. The measurements
in this paper are with Shirani only.}

\subsection{Front-tracking: explicit polyline}

In the plane the front is a closed, ordered polyline of markers with explicit
connectivity, advected by the interpolated velocity and maintained by local
surgery --- splitting an edge that grows beyond a tolerance, merging one that
falls below it. Surface tension enters as a force distributed from the markers
to the grid.

\caveat{The two-dimensional front-tracking core carries a header comment
describing it as phase B1 --- kinematic advection in a prescribed field,
without the solver. The coupling is supplied by a separate adapter in
\texttt{examples-common}, and the example's build file describes the assembled
pair as phase~B2. Before any production claim is made from this path, the
degree of coupling actually exercised has to be stated precisely; the 301-frame
Newtonian run on disk indicates a coupled run, but that is an inference from
the output, not a reading of the call graph.}

\subsection{Interface extraction for comparison}

Both methods are measured through the \emph{same} quantities, computed by the
same code path: the integrals that define circularity, centre of mass and rise
velocity are taken over the volume-fraction field, which VOF advects directly
and which front-tracking writes from the geometry of the front. The measurement
is therefore identical on both sides and any difference that survives belongs
to the method.

\subsection{Adaptive refinement}

\missing{the two-dimensional criterion, as implemented: seeds from the
interface, dilation to the refinement levels, and the minimum-cell rule.}

The one result already established here is negative and worth stating early,
because it governs how any convergence study in this paper must be set up. With
the refined band specified as a fixed \emph{number of cells} per side, each
resolution refines a \emph{different physical width} --- the finest mesh
refines the largest fraction of the domain --- and a mesh-convergence study
built that way does not converge for reasons that have nothing to do with the
discretisation.

\missing{the measured evidence for the paragraph above: five series, none
monotone, with non-monotonicity between 12 and 40 times the sampling floor.
Port the numbers from the report and put them in \texttt{dados/}.}

The remedy is to anchor the band to a fixed \emph{physical} width, taken from
the coarsest resolution of the series, so that finer meshes only gain cells.
A floor on the number of fine cells per side is kept and \emph{overrides} the
physical band when the two disagree, because the reconstruction and the
fraction transport both need that neighbourhood: it is a correctness
constraint, not a tuning parameter.

\section{Verification}

\status{Partial. The two-dimensional oracles have not been inventoried with
the discipline applied to the three-dimensional set, where twelve harnesses
each carry an exit code and a \emph{vacuity guard} --- they count how many
cases they actually exercised and fail if that count is zero. A harness that
never touches the case under test prints ``0 failures'', which is
indistinguishable from passing. Until the 2D harnesses carry the same guard,
their green cannot be quoted.}

\missing{the inventory of two-dimensional oracles, each with the quantity it
measures, the closed form it is measured against, and the measured error.
Candidates already present in the tree: the Laplace pressure jump
$\Delta p = \sigma/R$ for a static drop; the LeVeque single-vortex advection
reversal; curvature against circle and ellipse; and the partition-of-unity
check on the volume fraction.}

\subsection{What the green does not prove}

\missing{the same caution as the 3D companion. Oracles of closed form test the
components in isolation; none of them exercises the coupled transient that the
production cases run. A passing oracle set bounds the component error, not the
solution error.}

\section{Newtonian two-dimensional bubble}\label{sec:newtonian}

\status{Measured and closed. Both methods ran the same case to $t=3$ ---
12\,000 steps, no restarts --- and the five benchmark targets are compared
against the published reference below. The central finding of the section is
not any one of the five but their \emph{sign}.}

\subsection{Case definition}

The reference case is the rising-bubble benchmark of Hysing
et~al.~\missing{cite}, which fixes the geometry, the fluid properties and the
boundary conditions and --- unusually --- publishes reference series from three
independent codes, so that the disagreement \emph{among} the references is
itself quantified. That spread is the resolution limit of the oracle: asking a
new code to do better than it would be asking for more accuracy than the
reference possesses.

\begin{adjustbox}{max width=\linewidth}
\pgfplotstabletypeset[
  string type,
  columns/quantity/.style={column name={quantity}},
  columns/TP2D/.style={column name={TP2D}},
  columns/FreeLIFE/.style={column name={FreeLIFE}},
  columns/MooNMD/.style={column name={MooNMD}},
  columns/consensus/.style={column name={consensus}},
  every head row/.style={before row=\toprule, after row=\midrule},
  every last row/.style={after row=\bottomrule},
]{\dat{hysing-referencia.dat}}
\end{adjustbox}

\subsection{The three benchmark quantities}

The benchmark is defined by three series, computed here exactly as its
Section~2.5 defines them:
%
\begin{align}
  c  &= \frac{2\sqrt{\pi A}}{P}, &
  y_c &= \frac{\int_{\text{bubble}} y \,\mathrm{d}x}{\int_{\text{bubble}} \mathrm{d}x}, &
  V_c &= \frac{\int_{\text{bubble}} u_y \,\mathrm{d}x}{\int_{\text{bubble}} \mathrm{d}x} .
\end{align}

\caveat{The circularity above is the perimeter of the area-equivalent circle
over the perimeter of the bubble. The form $4\pi A/P^2$, common in
implementations, is its \emph{square}. We used the quadratic form at one point
and compared it with the published value without converting, reporting $0.7987$
against a published $0.9013$ when the comparable number was
$\sqrt{0.7987} = 0.8937$. The discrepancy looked like a method error and was a
definition error. Both forms are called ``circularity'' in the literature.}

Both integrals are taken over the volume-fraction field --- which VOF advects
and which front-tracking writes from the geometry of the front --- so the
measurement is the same on both sides and any difference that survives belongs
to the method.

\subsection{The runs}

\begin{adjustbox}{max width=\linewidth}
\pgfplotstabletypeset[
  string type,
  columns/quantity/.style={column name={quantity}},
  columns/FT/.style={column name={front-tracking}},
  columns/VOF/.style={column name={VOF}},
  every head row/.style={before row=\toprule, after row=\midrule},
  every last row/.style={after row=\bottomrule},
]{\dat{hysing-final.dat}}
\end{adjustbox}

\subsection{The two methods agree with each other, until they do not}

Early in the run the agreement between the two methods is far closer than the
agreement of either with the reference: the centre-of-mass series differ in the
fifth decimal while both sit a few tenths of a percent below the published
values.

\begin{adjustbox}{max width=\linewidth}
\pgfplotstabletypeset[
  string type,
  columns/t/.style={column name={$t$}},
  columns/ycft/.style={column name={$y_c$ FT}},
  columns/ycvof/.style={column name={$y_c$ VOF}},
  columns/difyc/.style={column name={$\Delta y_c$}},
  columns/vcft/.style={column name={$V_c$ FT}},
  columns/vcvof/.style={column name={$V_c$ VOF}},
  columns/difvc/.style={column name={$\Delta V_c$}},
  columns/circft/.style={column name={$c$ FT}},
  every head row/.style={before row=\toprule, after row=\midrule},
  every last row/.style={after row=\bottomrule},
]{\dat{hysing-parcial.dat}}
\end{adjustbox}
\captionof{table}{Both methods over the early window.}\label{tab:parcial}


\begin{figure}[htbp]
\centering
\begin{adjustbox}{max width=\linewidth}
\begin{tikzpicture}
\begin{axis}[
  name=vc, width=7.2cm, height=5.0cm,
  xlabel={$t$}, ylabel={$V_c$}, grid=major,
  title={rise velocity, early window}, title style={font=\small},
  legend style={font=\scriptsize, at={(0.5,-0.24)}, anchor=north},
  legend columns=2, legend cell align=left,
]
\addplot[thick,mark=*,mark size=1.1]          table[x=t,y=vcft]  {\dat{hysing-parcial.dat}};
\addlegendentry{front-tracking}
\addplot[densely dashed,mark=o,mark size=1.4] table[x=t,y=vcvof] {\dat{hysing-parcial.dat}};
\addlegendentry{VOF}
\end{axis}
\begin{axis}[
  at={(vc.right of south east)}, anchor=left of south west,
  width=7.2cm, height=5.0cm,
  xlabel={$t$}, ylabel={$\Delta V_c$}, grid=major,
  ymode=log, ymin=5e-5, ymax=2e-2,
  title={divergence between the methods}, title style={font=\small},
]
\addplot[thick,mark=*,mark size=1.1] table[x=t,y=difvc] {\dat{hysing-divergencia-plot.dat}};
\end{axis}
\end{tikzpicture}
\end{adjustbox}
\caption{Left: the rise velocity of both methods over the early window, where
the two are indistinguishable --- the curves overlap to plotting accuracy and
the tabulated difference is at most \num{7.4e-5} (Table~\ref{tab:parcial}).
Right: the difference between them over the window in which it grows, on a
logarithmic axis; it grows by a factor of $68$ --- from \num{1.1e-4} to
\num{7.5e-3} --- between $t=0.4$ and $t=1.6$, which is the deformation
phase, and it is not monotone: the dip at $t=1.4$ is in the stored series. \textbf{Both panels are sampled
windows, not the full run}: the stored series cover $t \le 0.35$ and
$0.4 \le t \le 1.6$ respectively.
\missing{the continuous series over the whole run to $t=3$. It is recoverable
without new simulation --- both runs left 301 frames on disk --- but it is a
re-reduction, not a plot of stored numbers.}}
\label{fig:hysing-series}
\end{figure}

That mutual agreement does not survive the deformation. By $t=1.6$ the two
methods differ by $3.4\%$ in $V_c$, and then --- this is the part worth
noticing --- they come back together in $y_c$: $1.078587$ against $1.078645$ at
$t=3$, a difference of \num{5.8e-5}.

\begin{adjustbox}{max width=\linewidth}
\pgfplotstabletypeset[
  string type,
  columns/t/.style={column name={$t$}},
  columns/difyc/.style={column name={$\Delta y_c$}},
  columns/difvc/.style={column name={$\Delta V_c$}},
  columns/difvcrelativa/.style={column name={$\Delta V_c$ relative}},
  every head row/.style={before row=\toprule, after row=\midrule},
  every last row/.style={after row=\bottomrule},
]{\dat{hysing-divergencia.dat}}
\end{adjustbox}
\captionof{table}{Divergence between the two methods.}\label{tab:divergencia}

This is the integral-versus-derivative asymmetry: $y_c$ smooths what $V_c$
amplifies. A comparison made on centres of mass alone would conclude excellent
agreement and miss the entire discrepancy --- which is an argument for
reporting the derivative quantity even when it is the less flattering one.

\subsection{Closure: the five targets}

\begin{adjustbox}{max width=\linewidth}
\pgfplotstabletypeset[
  string type,
  columns/target/.style={column name={target}},
  columns/method/.style={column name={method}},
  columns/ours/.style={column name={ours}},
  columns/published/.style={column name={published}},
  columns/deviation/.style={column name={deviation}},
  columns/inunitsofthespread/.style={column name={in units of the spread}},
  every head row/.style={before row=\toprule, after row=\midrule},
  every last row/.style={after row=\bottomrule},
]{\dat{hysing-fechamento.dat}}
\end{adjustbox}

The last column is the one that matters: it measures the deviation in units of
the disagreement the three published codes have \emph{among themselves}.
\textbf{Three of the five fall inside it.}

\subsection{The finding that outweighs the five numbers}

\textbf{All five deviations are negative.} None on the positive side. The
bubble comes out less fast, less circular and lower than the literature, in
\emph{two} independent discretisations --- one Lagrangian and one Eulerian,
with curvature and force computations that have nothing in common.

Five out of five in the same direction is not a plausible coincidence. It is
the signature of a \textbf{shared bias}: the defect lies in what the two
methods have in common --- mesh, time step, boundary conditions, momentum
discretisation --- and not in the interface treatment.

This also corrects a conclusion we had drawn from a single target. When only
$V_{c,\max}$ had closed, the gap between front-tracking ($0.31\times$ the
spread) and VOF ($2.29\times$) pointed at the interface treatment. With all
five on the table the reading splits, and the two halves are compatible: the
\emph{difference between the methods} does answer to the interface treatment,
and in the expected direction, with the curvature of the explicit geometry
beating the curvature of a smoothed field; the \emph{common bias}, which
neither avoids, does not.

\begin{adjustbox}{max width=\linewidth}
\pgfplotstabletypeset[
  string type,
  columns/quantity/.style={column name={quantity}},
  columns/FT/.style={column name={front-tracking}},
  columns/VOF/.style={column name={VOF}},
  columns/published/.style={column name={published}},
  every head row/.style={before row=\toprule, after row=\midrule},
  every last row/.style={after row=\bottomrule},
]{\dat{hysing-vcmax.dat}}
\end{adjustbox}

Three candidates for the bias, \textbf{none of them measured}: the single
resolution (40 cells per bubble radius), the capillary time-step restriction,
and the possibility that the wake refinement is insufficient. A convergence
study in $h$ separates the three, and it is the first item on the list of next
steps --- now with a measured reason to be first rather than for methodological
completeness.

\subsection{Shape}
\begin{figure}[htbp]
\centering
\begin{tikzpicture}
\begin{axis}[
  name=evol, width=6.2cm, height=9.2cm, axis equal image,
  xlabel={$x$}, ylabel={$y$}, grid=major,
  xmin=0.11, xmax=0.89, ymin=0.22, ymax=1.32,
  title={the bubble, $t=0$ to $t=3$}, title style={font=\small},
  legend style={font=\scriptsize, at={(0.5,-0.13)}, anchor=north},
  legend columns=2, legend cell align=left,
]
\addplot[thick]           table[x=x,y=y] {\dat{formas/hy_ft_0.dat}};
\addlegendentry{front-tracking}
\addplot[densely dashed]  table[x=x,y=y] {\dat{formas/hy_vof_0.dat}};
\addlegendentry{VOF}
\addplot[thick]           table[x=x,y=y] {\dat{formas/hy_ft_1.dat}};
\addplot[densely dashed]  table[x=x,y=y] {\dat{formas/hy_vof_1.dat}};
\addplot[thick]           table[x=x,y=y] {\dat{formas/hy_ft_2.dat}};
\addplot[densely dashed]  table[x=x,y=y] {\dat{formas/hy_vof_2.dat}};
\addplot[thick]           table[x=x,y=y] {\dat{formas/hy_ft_3.dat}};
\addplot[densely dashed]  table[x=x,y=y] {\dat{formas/hy_vof_3.dat}};
\end{axis}
\begin{axis}[
  at={(evol.right of south east)}, anchor=left of south west,
  width=6.6cm, height=6.6cm, axis equal image,
  xlabel={$x$}, grid=major,
  xmin=0.147, xmax=0.167, ymin=0.794, ymax=0.814,
  title={$t=1.9$, magnified 34 times}, title style={font=\small},
  scaled ticks=false, tick label style={font=\tiny, /pgf/number format/fixed,
                                        /pgf/number format/precision=3},
  xtick={0.150,0.155,0.160,0.165}, ytick={0.795,0.800,0.805,0.810},
]
\addplot[thick]           table[x=x,y=y] {\dat{formas/hy_ft_2.dat}};
\addplot[densely dashed]  table[x=x,y=y] {\dat{formas/hy_vof_2.dat}};
\end{axis}
\end{tikzpicture}
\caption{The interface in the two methods at $t=0$, $1$, $1.9$ and $3$. The
bubble rises from $y_c=0.5$ to $1.08$ and flattens: the width-to-height ratio
goes from $1.00$ to $1.87$, and $t=1.9$ is the instant of maximum deformation
--- the same instant as the minimum circularity. From $1.9$ to $3$ the shape
barely changes ($1.87$ to $1.86$): the bubble reaches a steady regime. The two
curves are \emph{indistinguishable} over the whole evolution; the right-hand
panel magnifies 34 times the point of \emph{greatest} discrepancy at the
instant of maximum deformation, where they separate by \num{4.7e-4} --- about
\SI{8}{\percent} of one fine cell. The VOF curve is the
$\mathrm{FracVol}=0.5$ contour by marching squares.}
\label{fig:hysing-formas}
\end{figure}


\begin{adjustbox}{max width=\linewidth}
\pgfplotstabletypeset[
  string type,
  columns/t/.style={column name={$t$}},
  columns/razaolarguraaltura/.style={column name={width/height}},
  columns/areaft/.style={column name={$A$ FT}},
  columns/areavof/.style={column name={$A$ VOF}},
  columns/perimft/.style={column name={$P$ FT}},
  columns/perimvof/.style={column name={$P$ VOF}},
  columns/cdosdois/.style={column name={$c$, both}},
  every head row/.style={before row=\toprule, after row=\midrule},
  every last row/.style={after row=\bottomrule},
]{\dat{hysing-forma.dat}}
\end{adjustbox}

\subsection{Under an adaptive mesh, VOF loses exact conservation}

On a fixed mesh the VOF fraction is conserved to machine precision. Under
adaptive refinement it is not, and the loss is not noise: the samples are
almost monotone.

\begin{adjustbox}{max width=\linewidth}
\pgfplotstabletypeset[
  string type,
  columns/method/.style={column name={method}},
  columns/fixed-mesh/.style={column name={fixed mesh}},
  columns/with-AMR-until-t/.style={column name={with AMR, until $t$}},
  columns/monotone-samples/.style={column name={monotone samples}},
  every head row/.style={before row=\toprule, after row=\midrule},
  every last row/.style={after row=\bottomrule},
]{\dat{hysing-massa-amr.dat}}
\end{adjustbox}

Front-tracking, which does not store the fraction, is affected far less. The
asymmetry is a property of the representation and not of the refinement
machinery, which is shared.

\subsection{Case 2: the skirted regime}

\missing{not run. Worth doing precisely because the published reference
solutions disagree with each other there: it is the regime that discriminates
between interface representations rather than confirming both.}

\section{Viscoelastic two-dimensional bubble}

\status{Not measured. This is the section that separates this paper from a
reproduction of the Newtonian benchmark, and it is also the one with the least
behind it. On the VOF side there are two short exploratory runs --- an
Oldroyd-B and a gPTT configuration, 101 steps each, 20 frames --- which are
enough to show the path executes and not enough to state anything. On the
front-tracking side there are viscoelastic input decks but no production run.
A matched pair has to be run before this section exists.}

\subsection{Case definition}

\missing{the viscoelastic case, chosen so that it is the Newtonian case of
Section~\ref{sec:newtonian} with one parameter changed. Keeping the geometry
and the density and viscosity ratios fixed is what makes the comparison
attributable to rheology; changing the case at the same time as the rheology
would confound the two.}

\subsection{What the Newtonian case cannot show}

\missing{the argument, once there are measurements to carry it. The
three-dimensional companion found that the discriminating feature between the
two interface representations was a \emph{shape} feature --- a trailing cusp
that a cell-wise planar reconstruction cannot resolve below the cell while an
explicit surface keeps resolving below it. Whether the planar analogue exists,
and whether it discriminates, is the question this section has to answer.}

\subsection{Uniform mesh}

\missing{the matched pair: front-tracking and VOF, same deck, same time step,
at two resolutions. Model and parameters to be fixed with the case definition.}

\subsection{Adaptive mesh}

\missing{depends on the uniform pair and on the adaptive criterion.}

\section{Cost and performance}

\status{Not measured in two dimensions. The three-dimensional companion
reports wall-clock rates per step on a uniform mesh and the cost of the
adaptive path against it; the same measurements have not been made here. The
runs that exist were not timed under controlled conditions.}

\missing{cost per step for both methods on the uniform mesh, at both
resolutions, with the machine and the core count stated; the cost of the
interface step as a fraction of the total; and the adaptive path against the
uniform one at matched accuracy rather than matched cell count --- the latter
comparison flatters refinement and says nothing.}

\section{Conclusions}

\status{Depends on everything above. Written last.}

\missing{conclusions}


% \nocite{*} faz a bibliografia do esqueleto listar TODAS as entradas
% disponiveis em referencias.bib, inclusive as ainda nao citadas no texto.
% REMOVER NA SUBMISSAO -- sem isso, uma secao sem \cite sai com bibliografia
% vazia e o PDF compila sem reclamar.
\nocite{*}

\bibliographystyle{elsarticle-num}
\bibliography{referencias}

\end{document}
